\documentclass{article}
\date{}

\usepackage[norsk]{babel}
\usepackage{subcaption}
\usepackage{amsmath}
\usepackage{graphicx}
\usepackage[a4paper, top=1in, bottom=1in, left=1in, right=1in]{geometry}
\usepackage{amsfonts}
\usepackage{dirtytalk}
\usepackage{enumerate}
\usepackage{amssymb}
\usepackage{float}
\usepackage{bm}


% vertical vector = \vect{...}
\newcommand{\vect}[1]{\begin{pmatrix} #1 \end{pmatrix}}
\newcommand{\vvect}[1]{\begin{vmatrix} #1 \end{vmatrix}}
\newcommand{\dett}[1]{\begin{vmatrix} #1 \end{vmatrix}}
% double underline is \uu{...}
\newcommand{\uu}[1]{\underline{\underline{#1}}}
% double bold underline is \uub{...}
\newcommand{\uub}[1]{\underline{\underline{\textbf{#1}}}}
% align counter reset
\AtBeginEnvironment{align}{\setcounter{equation}{0}}

% lra : <->
\newcommand{\lra}{\leftrightarrow}
% eq : <=>
\newcommand{\eq}{\quad\Longleftrightarrow\quad}
% req : =>
\newcommand{\req}{\quad\Longrightarrow\quad}
% imp : =>
\newcommand{\imp}{\quad\implies\quad}
% ra : ->
\newcommand{\ra}{\quad\rightarrow\quad}

% 3x space
\newcommand{\trip}{\ \ \ }
% relational equal
\newcommand{\eqto}[1]{\overset{#1}{\sim}}
\newcommand{\eqq}{\quad=\quad}
\newcommand{\eqqto}[1]{\overset{\substack{#1}}{\sim}}


\newcommand{\lrule}[1]{%
    \vspace{#1pt} % Space above the line
    \hrule % Draw the horizontal line
    \vspace{#1pt} % Space below the line
}

\makeatletter
\renewcommand*\env@matrix[1][*\c@MaxMatrixCols c]{%
  \hskip -\arraycolsep
  \let\@ifnextchar\new@ifnextchar
  \array{#1}}
\makeatother

\setlength{\parskip}{1em}
\setlength{\parindent}{0pt}

\begin{document}

\section*{Seksjon 2.1}

\subsection*{Oppgave 1}
Sann / usann

\begin{enumerate}[a)]

    \item \say{Anta polynomet $p_n$ interpolerer $f$ i $n+1$ punkter, og at polnyomet $q_n$ interpolerer $g$ i de samme punktene. Det interpolerende polynomet $f+g$ i de samme punktene er da lik $p_n+q_n$}
    
\end{enumerate}

\subsection*{Oppgave 4}
 Bevis for at interpolasjonspolynomet er unikt

 \begin{enumerate}[a)]
 
     \item Anta det finnes to kvadratiske polynomer $p_1$ og $p_2$ som interpolerer funksjonen $f$ i de tre punktene $x_0, x_1, x_2$. \\
     Se på differansen $p=p_2-p_1$ \\
     Hva er verdiene t il $p$ i interpolasjonspunktene?
         \\\\
    \(p(x) = p_2(x) - p_1(x)\)
    \\\\
    - Siden både $p_1$ og $p_2$ interpolerer på de samme punktene, vil både $x$ og $y$ verdiene være like. dvs\\
    \[
    \begin{array}{c  c}
    p_1(x_1) = y_1 & p_2(x_1) = y_1 \\
    p_1(x_2) = y_2 & p_2(x_2) = y_2 \\
    p_1(x_3) = y_3 & p_2(x_3) = y_3 \\
    \end{array}
    \]

    Ser vi på uttrykket $p$ igjen så får vi for hvert punkt:\\\\
    $p(x_1) = p_2(x_1) - p_1(x_1) \eqq y_1 - y_1 \eqq 0$ \\
    $p(x_2) = p_2(x_2) - p_1(x_2) \eqq y_2 - y_2 \eqq 0$ \\
    $p(x_3) = p_2(x_3) - p_1(x_3) \eqq y_3 - y_3 \eqq 0$ \\\\

    Altså alle verdiene er 0, dvs. \uu{$P(x_i)=0$, for $i \in [0,3]$}


     \vspace{2em}

     \item Bruk observasjonen i (a) til å vise at $p_1$ og $p_2$ må være like

     \vspace{2em}

     \item Generaliser resultatet fra (a) og (b) til polynomer av grad $n$
     
 \end{enumerate}


\newpage
\section*{Seksjon 2.2}

\subsection*{Oppgave 1}
\(f(x)=x^2\) med polynom av grad 3 i punktene $0, 1, 2, 3$

\(
\begin{array}{c | c | c | c}
     f(0) = 0 & c_0 = 0          &    c_0 = 0 & c_0 = 0\\
     f(1) = 1 & c_0 + c_1(1)   & c_1 = 1 & c_1 = 1\\
     f(2) = 4 & c_0 + c_1(2) + c_2(2)(2-1) &  c_2 = \frac{4 - 2}{2} & c_2 = 1\\
     f(3) = 9 & c_0 + c_1(3) + c_2(3)(3-1) + c_3(3)(3-1)(3-2) 
     & c_3 = \frac{9 - 3  - 6}{6} &  c_3 = 0
\end{array}
\)\\

For 
\(P_3(4):\\\\
 f(4) = 16 
 \\\\
 c_0 + c_1(4) + c_2(4)(4-1) + c_3(4)(4-1)(4-2) + c_4(4)(4-1)(4-2)(4-3)
 \\\\
 4+12+24c_4 = 16 \req 24c_4 = 16-4-12 \req c_4 = 0
\)\\


Newtonpolynomet med $P_3(4)$ blir:

\(p_3(x) = x + x(x-1)  = x^2\)\\

Det vil si at 
\(p_3(4) = 4 + 4(4-1) = 4 + 4(3) = 4 + 12 = 16\)

\newpage
\subsection*{Oppgave 2a}
Har blitt gitt data:

\(
\begin{array}{c | c | c | c | c}
    x & 0 & 1 & 3 & 4 \\ 
    \hline 
    f(x) & 1 & 0 & 2 & 1 \\ 
\end{array}
\)

\(p_3(x)=c_0(x-1)(x-3)(x-4)+c_1x(x-3)(x-4)+c_2x(x-1)(x-4)+c_3x(x-1)(x-3)
\)

Interpolasjonsbetingelsene sier at $P_3(x_i) = f(x_i)$, derfor:

\(p_3(0)=1 \req c_0(-1)(-3)(-4) = 1 \req -12c_0 = 1 \req c_0 = -\frac{1}{12}
\)

\(p_3(1) = 0 \req c_1(1)(-2)(-3) = 0 \req 6c_1 = 0 \req c_1 = 0
\)

\(p_3(3) = 2 \req c_2(3)(2)(-1) = 2 \req -6c_3 = 2 \req c_2 = -\frac{1}{3}
\)

\(p_3(4) = 1 \req c_3(4)(3)(1) = 1 \req 12c_3 = 1 \req c_3 = \frac{1}{12}
\)\\\\

Koeffisientene fra interpolasjonsbetingelsene:

\(\uu{c_0 = -\frac{1}{12}, \quad c_1 = 0, \quad c_2 = -\frac{1}{3}, \quad c_3 = \frac{1}{12}}\)\\\\

Dobbeltsjekker ved å sette inn:

\(p_3(x) = -\frac{1}{12}(x-1)(x-3)(x-4)-\frac{1}{3}(x)(x-1)(x-4) + \frac{1}{12}(x)(x-1)(x-3)\)

\(p_3(0) = -\frac{1}{12}(-1)(-3)(-4) =\frac{12}{12} = 1\)

\(p_3(1) = 0 \text{      (x-1) som blir 0 i hvert ledd, da blir hele 0}\)

\(p_3(3) = -\frac{1}{3}(3)(2)(-1) = \cdot\frac{6}{3} = 2\)

\(p_3(4) = \frac{1}{12}(4)(3)(1) = \frac{12}{12} = 1\)

\subsection*{Oppgave 2b}
Newtonformen til det interpolerende polynomet: bruker c-verdiene funnet i a)

\(p_3(x)=c_0 + c_1(x-x_0) + c_2(x-x_0)(x-x_1) + c_3(x-x_0)(x-x_1)(x-x_2)\)

Med interpolasjonsbetingelsene:

\(
\begin{array}{c | c | c | c | c}
     p_3(0) = 1 &   & &  & c_0 = 1\\
     p_3(1) = 0 &  c_0 + c_1 = 0 & c_1 = -c_0 &  & c_1 = -1 \\
     p_3(3) = 2 & c_0 + 3c_1 + 6c_2 = 2 & c_2 = \frac{2 - c_0 - 3c_1}{6} & c_2 = \frac{2-1+3}{6} & c_2 = \frac{2}{3} \\
     p_3(4) = 1 & c_0 + 4c_1 + 12c_2 + 12c_3 = 1 & c_3 = \frac{1 - c_0 - 4c_1 -12c_2}{12} & c_3 = \frac{1-1+4-8}{12} & c_3 = -\frac{1}{3}
\end{array}
\)

Det interpolerende polynom på newtonform er da:

\[\uu{p_3(x) = 1- x + \frac{2}{3}x(x-1) - \frac{1}{3}x(x-1)(x-3)}\]

\newpage

\subsection*{Oppgave 2c}
Verifisering av (a) og (b) er like:

$p_3L(x)$ er lagrangeformen og $p_3N(x)$ er newtonformen

\vspace{1em}

\(
p_3L(x) = -\frac{1}{12}(x-1)(x-3)(x-4)-\frac{1}{3}(x)(x-1)(x-4) + \frac{1}{12}(x)(x-1)(x-3)
\)
\[
p_3L(0) = -\frac{1}{12}(-1)(-3)(-4) = \frac{12}{12} = 1
\]
\[
p_3L(1) = 0, \text{  alle ledd inneholder } (x-1)
\]
\[
p_3L(3) = -\frac{1}{3}(3)(2)(-1) = \frac{6}{3} = 2
\]
\[
p_3L(4) = \frac{1}{12}(4)(3)(1) = \frac{12}{12} = 1
\]

\vspace{2em}

\(
p_3N(x) = 1- x + \frac{2}{3}x(x-1) - \frac{1}{3}x(x-1)(x-3)
\)
\[
p_3N(0) = 1-0 = 1
\]
\[p_3N(1) = 1-1 = 0\]
\[p_3N(3) = 1-3 + \frac{2}{3}(3)(2) = 1-3+4 = 2\]
\[p_3N(4) = 1 - 4 + \frac{2}{3}(4)(3) - \frac{1}{3}(4)(3) = 1 - 4 + 8 -4 = 1\]

Vi kan se at:

\(
p_3L(0) = p_3N(0) = 1 \\
p_3L(1) = p_3N(1) = 0 \\
p_3L(3) = p_3N(3) = 2 \\
p_3L(4) = p_3N(4) = 1
\)

Som viser at (a) og (b) er like.

\newpage

\subsection*{Oppgave 3a}
Har blitt gitt data:

\(
\begin{array}{c | c | c | c}
    x & 0 & 1 & 2\\ 
    \hline 
    f(x) & 2 & 1 & 0 \\ 
\end{array}
\)

Newtonformen til det kvadratiske interpolerende polynomet:

Interpoleringsbetingelsene:

\(p_3(0) = 2 \req c_0 = 2\)

\(p_3(1) = 1 \req c_0 + c_1(1) = 1 \req 2 + c_1 = 1 \req c_1 = -1\)

\(p_3(2) = 0 \req c_0 + 2c_1 + 2c_2 = 0 \req 2-2+2c_2 = 0 \req c_2 = 0\)

\vspace{1em}

Gir det interpolerende polynomet:

\(p_3(x) = 2 - 1(x-0) + 0(x-0)(x-1) \eqq 2-1(x) \eqq 2 - x\)

\vspace{1em}

\textbf{Konklusjon:} \uu{Det er ingen forskjell mellom $y = 2 - x$ og newtonpolynomet $p_3(x) = 2 - x$}

\vspace{2em}
\subsection*{Oppgave 3b}

\end{document}

